% Plano de Laplace: polos de un oscilador subcrítico y su interpretación.
\documentclass[tikz,border=4pt]{standalone}
\usepackage{resonancia}
\begin{document}
\begin{tikzpicture}[scale=1]
  \def\R{3.2}
  \def\b{52}   % ángulo beta (grados)
  \coordinate (O) at (0,0);
  \coordinate (P) at ({-\R*cos(\b)},{\R*sin(\b)});
  \coordinate (Pc) at ({-\R*cos(\b)},{-\R*sin(\b)});
  % ejes
  \draw[eje] (-4.2,0) -- (1.2,0) node[right, etiqueta] {$\mathrm{Re}\,s=\sigma$};
  \draw[eje] (0,-3.7) -- (0,3.8) node[above, etiqueta] {$\mathrm{Im}\,s=\omega$};
  % semicircunferencia de radio omega_n
  \draw[punteada] (0,3.2) arc (90:270:3.2);
  % radio, proyecciones
  \draw[trazo, color=rAzul] (O) -- (P) node[pos=0.55, above left, etiqueta, color=rAzul] {$\omega_n=|s|$};
  \draw[punteada] (P) -- ({-\R*cos(\b)},0);
  \draw[punteada] (P) -- (0,{\R*sin(\b)});
  \node[etiqueta, left] at (0,{\R*sin(\b)}) {$\omega_d$};
  \node[etiqueta, below] at ({-\R*cos(\b)/2},0) {$\sigma=-\zeta\omega_n$};
  % ángulo beta
  \draw[color=rVerde, trazo] (-1.0,0) arc (180:{180-\b}:1.0);
  \node[etiqueta, color=rVerde] at (-1.42,0.42) {$\beta$};
  \node[etiqueta, color=rVerde, anchor=north west] at (0.12,-0.12) {$\zeta=\cos\beta$};
  % polos
  \node[polo, scale=1.7] at (P) {$\times$};
  \node[polo, scale=1.7] at (Pc) {$\times$};
  \node[etiqueta, above left] at (P) {$s_1=\sigma+i\omega_d$};
  \node[etiqueta, below left] at (Pc) {$s_1^*$};
  % trayectoria con zeta
  \draw[flecha, color=rGris] ({-\R*cos(25)},{\R*sin(25)}) arc (155:200:3.2);
  \node[etiqueta, color=rGris, anchor=west] at (-3.25,-1.35) {$\zeta\uparrow$};
  \node[etiqueta, color=rGris, anchor=south west] at (0.1,3.2) {$\zeta=0$};
  \node[etiqueta, color=rGris, anchor=north east] at (-3.2,-0.05) {$\zeta=1$};
\end{tikzpicture}
\end{document}
